% !TeX program = lualatex
% !TeX encoding = utf8
% !TeX spellcheck = uk_UA
% !TeX root = ../ClassicalElectrodynamics.tex

%=========================================================
\chapter{Електродинаміка у просторі Мінковського}\label{\currfilebase}
%=========================================================

Рівняння Максвелла, встановлені задовго до створення СТВ, чудово узгоджуються з принципом відносності. Тому їх можна використати для знаходження зв'язку між напруженостями електричного та індукціями магнітного полів в різних системах відліку, що й буде зроблено далі. Щодо рівнянь руху зарядів, то вони вимагають істотної модифікації у порівнянні з ньютонівською теорією.

%% --------------------------------------------------------
\section{Чотиривимірна форма рівнянь Максвелла}
%% --------------------------------------------------------

Згідно з експериментом, рівняння Максвелла задовільно описують електродинамічні явища у класичній\footnote{Тут класичній~= неквантовій.} області. Головними фігурантами рівнянь Максвелла є електричні заряди та напруженості полів\footnote{Тобто, напруженість електричного поля та індукція магнітного.}. Як відомо, однією з головних властивостей заряду є його інваріантність: величина заряду ізольованого тіла не залежить від руху, вона однакова в усіх інерціальних системах. Тепер треба з'ясувати, як перетворюються напруженості полів. Для цього враховуватимемо, що за принципом відносності рівняння електромагнітного поля\index{Електромагнітне поле} зберігають свій вигляд в усіх інерціальних системах.

Перепишемо рівняння Максвелла у чотиривимірних позначеннях, більш зручних для використання в теорії відносності. Нехай $\rho$ --- густина електричного заряду деякого елемента середовища; $\vect{j}=\rho\vect{v}$ --- звичайний тривимірний вектор густини струму; $\vect{v}=d\vect{x}/dt$ --- тривимірна швидкість зарядів у точці, що розглядається. Знайдемо трансформаційні властивості величини
\begin{equation*}
    j^{\mu} = \rho\frac{ds}{dt}\frac{dx^{\mu}}{ds},
    \qquad
    \{j^{\mu}\} \equiv \{j^{0},\vect{j}\},\quad
    j^{0} = c\rho,\quad
    j^{i} = \rho\,dx^{i}/dt.
\end{equation*}

З формули для перетворення густини заряду~\eqref{eq:charge_density_transform}, яку легко узагальнити на випадок довільного напрямку швидкості, випливає, що величина $\rho\sqrt{1-\vect{v}^{2}/c^{2}} = \rho\,ds/dt$ є інваріантом --- це густина заряду у власній системі елемента середовища. Тоді $j^{\mu}$ має ті ж трансформаційні властивості, що й $dx^{\mu}/ds$, тобто це контраваріантний вектор. Його називають \textbf{чотиривектором густини струму\index{Чотиривектор!густини струму}}. Цей вектор введено для середовища з неперервним розподілом зарядів та швидкостей, але його можна означити і для точкових зарядів $q_{i}$ з траєкторіями $x^{\alpha}=x^{\alpha}_{(i)}$ (індекс $i=1,2,\ldots$ нумерує заряди):
\begin{equation*}
    j^{\mu} = c\sum_{i}q_{i}\int ds_{i}\,\delta^{(4)}(x - x_{(i)}(s_{i}))
    \frac{dx_{(i)}^{\mu}}{ds_{i}},
\end{equation*}
де інтегрування проводиться вздовж траєкторій зарядів, $\delta^{(4)}(x)$ --- чотиривимірна дельта-функція Дірака\footnote{За визначенням, для будь-якої неперервної функції $f$: $\int\delta^{(4)}(x-x_0)f(x)\,d^4x = f(x_0)$.}.

Для опису електромагнітного поля\index{Електромагнітне поле} введемо антисиметричну матрицю
\begin{equation*}
\left\| F_{\mu\nu} \right\| =
\begin{pmatrix}
0 & E_1 & E_2 & E_3 \\
-E_1 & 0 & -B_3 & B_2 \\
-E_2 & B_3 & 0 & -B_1 \\
-E_3 & -B_2 & B_1 & 0
\end{pmatrix},
\end{equation*}
де $F_{\mu\nu}=-F_{\nu\mu}$; $E_i$ --- компоненти напруженості електричного поля\index{Напруженість електричного поля}, $B_i$ --- індукції магнітного поля\index{Індукція магнітного поля} в декартових координатах $x^{i}$ ($i=1,2,3$). Цю матрицю називають \emph{тензором електромагнітного поля\index{Тензор електромагнітного поля}}. Компоненти цього тензора можна означити ще й так:
\begin{equation*}
    F_{0i} = -F_{i0} = E_i, \qquad F_{ij} = -\varepsilon_{ijk}B_k,
\end{equation*}
де в останній рівності береться сума по $k$, причому індекси $i$, $j$, $k$ приймають значення 1,~2,~3. Як буде видно далі, $F_{\mu\nu}$ дійсно є двічі коваріантним тензором\index{Тензор!коваріантний} відносно перетворень групи Лоренца та групи Пуанкаре; відповідно $F^{\mu\nu}=\eta^{\mu\alpha}\eta^{\nu\beta}F_{\alpha\beta}$ є двічі контраваріантним тензором.

\textbf{Рівняння Максвелла\index{Рівняння Максвелла}} в гаусовій системі одиниць\index{Гаусова система одиниць}, переписані за допомогою $F$, мають вид:
\begin{equation}\label{eq:Maxwell_4D_1}
    \partial_{\mu}F^{\mu\nu} = \frac{4\pi}{c}j^{\nu},
\end{equation}
\begin{equation}\label{eq:Maxwell_4D_2}
    \partial_{\alpha}F_{\beta\gamma}
    + \partial_{\beta}F_{\gamma\alpha}
    + \partial_{\gamma}F_{\alpha\beta} = 0,
\end{equation}
де $\partial_{\mu}f \equiv \partial f/\partial x^{\mu}$. Зауважимо, що \eqref{eq:Maxwell_4D_2} можна записати також у вигляді
\begin{equation*}
    \varepsilon^{\mu\alpha\beta\gamma}\partial_{\alpha}F_{\beta\gamma} = 0,
\end{equation*}
де $\varepsilon^{\alpha\beta\gamma\delta}$ --- абсолютно антисиметричний символ Леві-Чівіта\index{Символ Леві-Чівіта}.

З антисиметрії $F_{\mu\nu}$ та з \eqref{eq:Maxwell_4D_1} випливає \textbf{закон збереження заряду\index{Закон збереження заряду}}:
\begin{equation}\label{eq:charge_conservation}
    \partial_{\mu}j^{\mu} = 0,
\end{equation}
як необхідна умова існування розв'язку рівнянь електромагнітного поля. Оскільки $j^{\mu}$ є чотири-вектор, а $\partial_{\alpha}j^{\beta}$ --- тензор, згортка $\partial_{\mu}j^{\mu}$ є інваріантом. Рівняння~\eqref{eq:charge_conservation} є інваріантним при перетвореннях Лоренца.

%=========================================================
\begin{problem}
Встановити прямим обчисленням еквівалентність рівнянь Максвелла у тривимірній та чотиривимірній формі.

\textit{Відповідь}: Рівняння~\eqref{eq:Maxwell_4D_1} при $\nu=0$ еквівалентне рівнянню
\[
    \Div\vect{E} = 4\pi\rho,
\]
рівняння~\eqref{eq:Maxwell_4D_1} при $\nu=1,2,3$ еквівалентне рівнянню
\[
    \Rot\vect{B} = \frac{4\pi}{c}\vect{j}
    + \frac{1}{c}\frac{\partial\vect{E}}{\partial t},
\]
рівняння~\eqref{eq:Maxwell_4D_2} при $\alpha,\beta,\gamma=1,2,3$ еквівалентне рівнянню
\[
    \Div\vect{B} = 0,
\]
а інші три комбінації $\alpha,\beta,\gamma$, що містять індекс «0», відповідають рівнянню
\[
    \Rot\vect{E} = -\frac{1}{c}\frac{\partial\vect{B}}{\partial t}.
\]
\end{problem}

%=========================================================
\begin{problem}
Встановити тривимірний аналог рівняння~\eqref{eq:charge_conservation}, користуючись означенням $j^{\mu}$.

\textit{Відповідь}: Рівняння~\eqref{eq:charge_conservation} є іншим записом \emph{рівняння неперервності\index{Рівняння неперервності}}:
\[
    \Div\vect{j} + \frac{\partial\rho}{\partial t} = 0.
\]
\end{problem}

%% --------------------------------------------------------
\section{Трансформаційні властивості \texorpdfstring{$F_{\mu\nu}$}{Fmunu}}
%% --------------------------------------------------------

Досі ми не з'ясовували трансформаційні властивості величин $F_{\mu\nu}$, які є розв'язками рівнянь електродинаміки~\eqref{eq:Maxwell_4D_1}, \eqref{eq:Maxwell_4D_2}. За певних граничних умов розв'язок цих рівнянь визначений єдиним чином. Для подальшого досить розглянути якусь одну постановку задачі для~\eqref{eq:Maxwell_4D_1}, \eqref{eq:Maxwell_4D_2}, що забезпечує єдиність, зокрема, що відповідає розгляду ізольованої системи. Адже фізичні властивості поля, в тому числі і трансформаційні співвідношення, що пов'язують різні інерціальні системи, не залежать від того, яка конфігурація зарядів створює це поле.

Врахуємо, що, за відсутності зовнішнього випромінювання, рівняння~\eqref{eq:Maxwell_4D_1}, \eqref{eq:Maxwell_4D_2} визначають електромагнітне поле $F_{\alpha\beta}$ (а також і $F^{\alpha\beta}$) обмеженої системи зарядів і струмів \emph{однозначно в усьому просторі}, якщо ці рівняння розглядають, починаючи з нескінченного минулого. Відповідний розв'язок для поля ізольованої системи можна подати за допомогою запізнюючих потенціалів (див. далі \eqref{eq:retarded_potentials} та представлення~\eqref{eq:4potential_def} електромагнітного поля через 4-вектор потенціалу), за допомогою яких $F_{\alpha\beta}$ \emph{однозначно} виражається через 4-вектор струму. Ці формули дають змогу явно дослідити трансформаційні властивості $F_{\alpha\beta}$ при переході в іншу інерціальну систему відліку. Але тут ми вчинимо інакше, виходячи безпосередньо з рівнянь Максвелла.

Нехай $F_{\alpha\beta}$, що подає напруженості полів, є розв'язком рівнянь Максвелла в інерціальній системі $S$ з координатами $\{x\}$, а $\tilde{F}_{\alpha\beta}$ --- розв'язком в інерціальній системі $S'$ з координатами $\{x'\}$: $x'^{\alpha'} = L^{\alpha'}{}_{\alpha}\,x^{\alpha}$.

Маючи компоненти $F_{\alpha\beta}$ та $F^{\alpha\beta}$ в системі $S$, можна формально ввести тензор в усіх системах координат відомими співвідношеннями, зокрема, в системі $S'$ це будуть компоненти
\begin{equation}\label{eq:Lorentz_transform_F}
    F'_{\mu\nu} = (L^{-1})^{\alpha}{}_{\mu}\,(L^{-1})^{\beta}{}_{\nu}\,F_{\alpha\beta},
    \qquad
    F'^{\mu\nu} = \eta^{\mu\alpha}\eta^{\nu\beta}F'_{\alpha\beta}
    = L^{\mu}{}_{\alpha}\,L^{\nu}{}_{\beta}\,F^{\alpha\beta}.
\end{equation}
За властивостями тензорів, ці величини можна диференціювати, утворюючи нові тензори, утворювати лінійні комбінації з тензорами аналогічної будови тощо. Оскільки, за припущенням, $F_{\alpha\beta}$ є розв'язком рівнянь Максвелла~\eqref{eq:Maxwell_4D_1}, \eqref{eq:Maxwell_4D_2} в системі $S$:
\begin{equation*}
    \frac{\partial F^{\alpha\nu}}{\partial x^{\nu}} + \frac{4\pi}{c}j^{\alpha} = 0,
    \qquad
    \partial_{\mu}F_{\nu\lambda} + \partial_{\nu}F_{\lambda\mu}
    + \partial_{\lambda}F_{\mu\nu} = 0,
\end{equation*}
аналогічне співвідношення маємо для тензорів в лівій частині рівнянь і в системі $S'$:
\begin{align*}
    \frac{\partial F'^{\mu\nu}}{\partial x'^{\nu}} + \frac{4\pi}{c}J'^{\mu}
    &= L^{\mu}{}_{\alpha}\!\left(\frac{\partial F^{\alpha\nu}}{\partial x^{\nu}}
    + \frac{4\pi}{c}j^{\alpha}\right) = 0, \\
    \partial'_{\alpha}F'_{\beta\gamma} + \partial'_{\beta}F'_{\gamma\alpha}
    + \partial'_{\gamma}F'_{\alpha\beta}
    &= (L^{-1})^{\mu}{}_{\alpha}(L^{-1})^{\nu}{}_{\beta}(L^{-1})^{\lambda}{}_{\gamma}
    \!\left(\partial_{\mu}F_{\nu\lambda} + \partial_{\nu}F_{\lambda\mu}
    + \partial_{\lambda}F_{\mu\nu}\right) = 0.
\end{align*}
Таким чином, тензор $F'_{\alpha\beta}$, що подається співвідношеннями~\eqref{eq:Lorentz_transform_F}, є розв'язком рівнянь електромагнітного поля в системі $S'$. Завдяки єдиності розв'язку цих рівнянь компоненти тензора $F'_{\alpha\beta} = \tilde{F}_{\alpha\beta}$ подають електромагнітне поле в новій системі відліку $S'$. Таким чином, зв'язок компонент тензора електромагнітного поля в різних системах дійсно можна подати формулами~\eqref{eq:Lorentz_transform_F}. Співвідношення, аналогічні~\eqref{eq:Lorentz_transform_F}, можна отримати між будь-якими системами відліку. Сукупність компонент $F_{\alpha\beta}$ в усіх системах утворює двічі коваріантний тензор.

%=========================================================
\begin{problem}\label{prb:3.3}
Відносно інерціальної системи відліку $S$ з координатами $\{t,\vect{r}\}$, $\vect{r}=(x,y,z)$, рухається із сталою швидкістю $\vect{v}$ інша інерціальна система $S'$ з координатами $\{t',\vect{r}'\}$, $\vect{r}'=(x',y',z')$. У деякий момент початки координат збігалися. Перетворення координат $S\to S'$ мають вид
\begin{equation*}
    \vect{r}'_{\|} = \frac{\vect{r}_{\|} - \vect{v}t}{\sqrt{1-\beta^2}},
    \qquad
    t' = \frac{t - (\vect{v}\cdot\vect{r})/c^2}{\sqrt{1-\beta^2}},
    \qquad
    \vect{r}'_{\bot} = \vect{r}_{\bot},
    \tag{$*$}
\end{equation*}
де позначки $\|$ та $\bot$ відповідають поздовжнім та поперечним компонентам тривимірних векторів. За допомогою закону перетворення для тензора електромагнітного поля покажіть, що відповідні до $(*)$ перетворення поперечних та поздовжніх компонент напруженості електричного поля та індукції магнітного поля мають вид:
\begin{align*}
    \vect{E}'_{\|} &= \vect{E}_{\|},
    \qquad
    \vect{B}'_{\|} = \vect{B}_{\|}, \\
    \vect{E}'_{\bot} &= \frac{\vect{E}_{\bot}
    + \tfrac{1}{c}[\vect{v}\times\vect{B}]}{\sqrt{1-\beta^2}},
    \qquad
    \vect{B}'_{\bot} = \frac{\vect{B}_{\bot}
    - \tfrac{1}{c}[\vect{v}\times\vect{E}]}{\sqrt{1-\beta^2}}.
\end{align*}
\end{problem}

%=========================================================
\begin{problem}
За допомогою формул перетворення електричного та магнітного полів при переході $S\to S'$ з вправи~\ref{prb:3.3} покажіть прямим обчисленням, що величини $\vect{E}^{2}-\vect{B}^{2}$ та $\vect{E}\cdot\vect{B}$ при цьому переході залишаються незмінними.
\end{problem}

%=========================================================
\begin{problem}
Виходячи з 4-вектора густини струму, покажіть, що перетворення густини заряду та густини струму, відповідні переходу $S\to S'$ з вправи~\ref{prb:3.3}, мають вид:
\[
    \rho' = \frac{\rho - (\vect{j}\cdot\vect{v})/c^2}{\sqrt{1-\beta^2}},
    \qquad
    \vect{j}'_{\|} = \frac{\vect{j}_{\|} - \rho\vect{v}}{\sqrt{1-\beta^2}},
    \qquad
    \vect{j}'_{\bot} = \vect{j}_{\bot}.
\]
\end{problem}

%% --------------------------------------------------------
\section{Рівняння руху зарядженої частинки}
%% --------------------------------------------------------

Рівняння Максвелла задовольняють принципу відносності; тому при створенні СТВ вони не потребували жодних коректив. Навпаки, рівняння руху ньютонівської механіки треба було модифікувати: вони не задовольняють принципам СТВ і призводять до нефізичних наслідків, особливо, коли швидкості руху наближаються до швидкості світла. Цю модифікацію можна отримати, якщо врахувати, що рівняння руху не містять якихось додаткових фізичних величин, окрім швидкості, прискорення, компонент $F_{\mu\nu}$, а також маси частинки $m$ та її заряду $q$.

Отримаємо рівняння руху зарядженої частинки з урахуванням \emph{принципу відповідності з класичною теорією}. У контексті рівнянь руху це означає, що у системі відліку, де швидкість частинки значно менша за швидкість світла $c$, можна користуватися ньютонівськими рівняннями
\begin{equation*}
    m\frac{d\vect{v}}{dt} = q\!\left\{\vect{E}
    + \frac{1}{c}[\vect{v}\times\vect{B}]\right\}.
\end{equation*}
Ми далі використаємо навіть більш слабку умову, що класичні рівняння руху виконуються при $\vect{v}=0$. Основна ідея полягає у тому, що співвідношення між тензорними величинами, записані в одній системі відліку, можна поширити на будь-які системи.

Розглянемо коваріантний 4-вектор
\begin{equation*}
    R^{\mu} = mc^2\frac{du^{\mu}}{ds} - qF^{\mu}{}_{\nu}\,u^{\nu},
\end{equation*}
де $s=c\tau$, $\tau$ --- власний час вздовж світової лінії частинки, $u^{\mu}=dx^{\mu}/ds$.

Зафіксуємо деякий момент $t_1$ і подію~А з координатами $(t_1,\vect{x}(t_1))$ та розглянемо інерціальну систему відліку $S_1$, таку, де швидкість частинки у цей момент дорівнює нулю: $\vect{v}=0$. Тоді, згідно з принципом відповідності, в цій системі у момент події~А можна використати ньютонівські співвідношення (в декартових координатах):
\begin{equation*}
    m\frac{d\vect{v}}{dt} = q\!\left\{\vect{E}
    + \frac{1}{c}[\vect{v}\times\vect{B}]\right\} = q\vect{E},
    \qquad
    m\frac{d(\vect{v}^2/2)}{dt} = 0.
\end{equation*}
Завдяки цьому легко обчислити усі компоненти в $S_1$: $R^{\mu}=0$ в момент події~А. Але, оскільки, за побудовою, $R^{\mu}$ це 4-вектор, згідно до властивостей векторів він дорівнює нулю у довільній інерціальній системі. Звідси маємо рівняння в момент події~А:
\begin{equation}\label{eq:rel_motion}
    mc^2\frac{du^{\mu}}{ds} = qF^{\mu}{}_{\nu}\,u^{\nu}.
\end{equation}
Ці міркування можна повторити для будь-якої події на траєкторії частинки, тому \textbf{рівняння~\eqref{eq:rel_motion} є шуканим релятивістським узагальненням рівняння руху зарядженої частинки\index{Рівняння руху!релятивістське}}.

Звернемо увагу, що серед чотирьох рівнянь~\eqref{eq:rel_motion} лише три є незалежними, якщо врахувати зв'язок між компонентами чотири-швидкості. Навпаки, можна показати, що вони не суперечать цьому зв'язку: якщо помножити~\eqref{eq:rel_motion} на $u_{\mu}$, в силу антисиметрії $F^{\mu\nu}$ у правій частині маємо $F^{\mu}{}_{\nu}\,u^{\nu}u_{\mu} \equiv F^{\mu\nu}u_{\mu}u_{\nu} \equiv 0$, звідки $u_{\mu}\,du^{\mu}/ds = \frac{1}{2}\,d(u_{\mu}u^{\mu})/ds = 0$, тобто з рівнянь~\eqref{eq:rel_motion} випливає $u^{\mu}u_{\mu}=\text{const}$. Кількість ступенів вільності, що визначається кількістю незалежних початкових умов для рівнянь руху, тут така ж сама, як і в класичній механіці.

З рівнянь руху отримаємо формулу для енергії та імпульсу частинки. Запишемо тривимірну форму рівнянь руху, яка є еквівалентною трьом компонентам~\eqref{eq:rel_motion} при $\mu=i=1,2,3$:
\begin{equation*}
    m\frac{d}{dt}\!\left(\frac{\vect{v}}{\sqrt{1-\vect{v}^2/c^2}}\right)
    = q\!\left\{\vect{E} + \frac{1}{c}[\vect{v}\times\vect{B}]\right\}.
\end{equation*}
Права частина --- це сила Лоренца\index{Сила Лоренца}, що визначає швидкість передачі імпульсу від поля до зарядженої частинки. Звідси випливає \emph{вираз для тривимірної частини релятивістського імпульсу}:
\begin{equation*}
    \vect{p} = \frac{m\vect{v}}{\sqrt{1-\vect{v}^2/c^2}}.
\end{equation*}
У згоді з принципом відповідності при $|\vect{v}|/c\ll 1$ маємо $\vect{p}\approx m\vect{v}$.

Компонента $\mu=0$ рівняння~\eqref{eq:rel_motion} дає
\begin{equation*}
    mc^2\frac{d}{dt}\!\left(\frac{1}{\sqrt{1-\vect{v}^2/c^2}}\right)
    = q(\vect{E}\cdot\vect{v}).
\end{equation*}
Права частина --- це звичайна потужність, що описує передачу енергії від поля до частинки за одиницю часу. Звідси випливає \emph{релятивістський вираз для повної енергії рухомої частинки}:
\begin{equation*}
    \mathcal{E} = \frac{mc^2}{\sqrt{1-\vect{v}^2/c^2}}.
\end{equation*}
Тут зроблено вибір константи інтегрування, що відповідає експериментальним даним щодо власної енергії частинок. Кінетична енергія частинки, за означенням, є $\mathcal{E}_k = \dfrac{mc^2}{\sqrt{1-\vect{v}^2/c^2}} - mc^2$.

%=========================================================
\begin{problem}
Частинка з зарядом $q$ та масою $m$ рухається в однорідному сталому електромагнітному полі, причому електричне $\vect{E}$ та магнітне $\vect{B}$ поля паралельні. В початковий момент $t=0$ частинка була у початку координат та мала швидкість $v$, перпендикулярну до полів. Знайти залежності координат зарядженої частинки від власного часу $\tau$.

\textit{Відповідь}. Нехай вісь $OZ$ напрямлена вздовж вектора напруженості електричного поля, а вісь $OX$ --- вздовж напрямку швидкості в початковий момент. Тоді
\begin{align*}
    ct &= \frac{c\gamma}{\omega_E}\sh(\omega_E\tau),
    &z &= \frac{c\gamma}{\omega_E}\!\left(\ch(\omega_E\tau)-1\right), \\
    x &= \frac{v\gamma}{\omega_H}\sin(\omega_H\tau),
    &y &= \frac{v\gamma}{\omega_H}\!\left(\cos(\omega_H\tau)-1\right),
\end{align*}
де $\omega_E = \dfrac{qE}{mc}$, $\omega_H = \dfrac{qH}{mc}$, $\gamma = \dfrac{1}{\sqrt{1-v^2/c^2}}$.
\end{problem}

%% --------------------------------------------------------
\section{Плоскі монохроматичні хвилі}
%% --------------------------------------------------------

За відсутності зарядів і струмів ($j^{\mu}=0$) розв'язок однорідних рівнянь Максвелла
\begin{equation*}
    \partial_{\nu}F^{\mu\nu} = 0, \qquad
    \partial_{\alpha}F_{\beta\gamma}
    + \partial_{\beta}F_{\gamma\alpha}
    + \partial_{\gamma}F_{\alpha\beta} = 0
\end{equation*}
можна шукати у вигляді суперпозиції плоских монохроматичних хвиль
\begin{equation*}
    F_{\mu\nu}(x) = f_{\mu\nu}\exp(-i\Omega),
\end{equation*}
де компоненти тензора $f_{\mu\nu}$ --- сталі, а скаляр $\Omega$ --- лінійна функція координат, $\Omega(x)=k_{\mu}x^{\mu}\equiv k^0x^0 - \vect{k}\vect{x}$; $k_{\mu}=\partial_{\mu}\Omega$ --- хвильовий (коваріантний) чотиривектор\index{Хвильовий вектор}.

Частота коливань\index{Ефект Доплера} періодичного процесу, яка сприймається спостерігачем на світовій лінії $x_{\text{obs}}^{\mu}(\tau)$, --- це швидкість зміни фази~$\Omega\!\left(x_{\text{obs}}^{\mu}(\tau)\right)$ за власним часом спостерігача:
\begin{equation}\label{eq:freq_observer}
    \omega = \frac{d}{d\tau}\,\Omega\!\left(x_{\text{obs}}^{\mu}(\tau)\right)
    = c\,\frac{d}{ds}\,\Omega = c\,k_{\mu}\,u_{\text{obs}}^{\mu},
\end{equation}
де $u_{\text{obs}}^{\mu}=dx_{\text{obs}}^{\mu}/ds_{\text{obs}}$ --- чотиришвидкість спостерігача, $\tau$ --- власний час спостерігача, що вимірює частоту, $s=c\tau=s_{\text{obs}}$.

Вектор $K^{\mu}=k^{\mu}-k_{\alpha}u_{\text{obs}}^{\alpha}\,u_{\text{obs}}^{\mu}$, ортогональний до $u_{\text{obs}}^{\mu}$, визначає напрямок руху електромагнітних хвиль відносно спостерігача. Дійсно, у власній системі спостерігача чотиришвидкість має лише одну ненульову компоненту $\{u_{\text{obs}}^{\mu}\}=\{1,0,0,0\}$, звідки $K^0=0$, а тривимірні компоненти векторів збігаються з $K^i=k^i$.

Підстановка тензора електромагнітного поля для плоских хвиль у \eqref{eq:Maxwell_4D_1} та \eqref{eq:Maxwell_4D_2} дає
\begin{equation}\label{eq:transversality}
    f^{\mu\nu}k_{\nu} = 0,
\end{equation}
\begin{equation}\label{eq:Bianchi_plane}
    f_{\alpha\beta}k_{\gamma} + f_{\beta\gamma}k_{\alpha} + f_{\gamma\alpha}k_{\beta} = 0.
\end{equation}
Перше рівняння~\eqref{eq:transversality} складає умову поперечності електромагнітних хвиль у вакуумі. Згортаючи~\eqref{eq:Bianchi_plane} з $k^{\gamma}$, після підсумовування за індексом $\gamma$ з урахуванням~\eqref{eq:transversality} отримуємо $f_{\alpha\beta}k_{\gamma}k^{\gamma}=0$ та, для ненульової амплітуди $f_{\alpha\beta}$,
\begin{equation}\label{eq:lightlike}
    k_{\gamma}k^{\gamma} = 0,
    \qquad\text{або}\qquad
    \eta^{\mu\nu}\frac{\partial\Omega}{\partial x^{\nu}}\frac{\partial\Omega}{\partial x^{\mu}} = 0.
\end{equation}
Звідси випливає $k^0=|\vect{k}|$.

Згортаючи~\eqref{eq:Bianchi_plane} з $f^{\alpha\beta}$ із врахуванням~\eqref{eq:transversality} дістанемо $f^{\alpha\beta}f_{\alpha\beta}=0$, або $F^{\alpha\beta}F_{\alpha\beta}=0$, що еквівалентне $\vect{E}^2=\vect{B}^2$. Згорнемо~\eqref{eq:Bianchi_plane} з $\varepsilon^{\mu\nu\alpha\beta}f_{\mu\nu}$ та врахуємо, що~\eqref{eq:Bianchi_plane} можна записати як $\varepsilon^{\mu\alpha\beta\gamma}k_{\alpha}f_{\beta\gamma}=0$. Після такої операції зникають доданки
\begin{equation*}
    \varepsilon^{\mu\nu\alpha\beta}(f_{\mu\nu}f_{\beta\gamma}k_{\alpha}
    + f_{\gamma\alpha}k_{\beta})
    \equiv (\varepsilon^{\mu\nu\alpha\beta}f_{\mu\nu}k_{\alpha})f_{\beta\gamma}
    + (\varepsilon^{\mu\nu\alpha\beta}f_{\mu\nu}k_{\beta})f_{\gamma\alpha} = 0
\end{equation*}
і залишається $\varepsilon^{\mu\nu\alpha\beta}f_{\mu\nu}f_{\alpha\beta}=0$, або $\varepsilon^{\mu\nu\alpha\beta}F_{\mu\nu}F_{\alpha\beta}=0$, що еквівалентне $\vect{E}\cdot\vect{B}=0$. Як і слід було очікувати, отримано звичайні співвідношення для плоских хвиль у вакуумі.

%% --------------------------------------------------------
\section{Ефект Доплера}
%% --------------------------------------------------------

Отримаємо формулу ефекту Доплера\index{Ефект Доплера}, який полягає у зміні спостережуваної частоти випромінювання в залежності від руху джерела відносно спостерігача. Аналогічно~\eqref{eq:freq_observer} власна частота коливань випромінювача з траєкторією $x_s^{\mu}(\tau_s)$, власним часом $\tau_s$ та 4-швидкістю $u_s^{\mu}=dx_s^{\mu}/ds_s$ є $\omega_0 = (d/d\tau_s)\,\Omega(x_s^{\mu}(\tau_s)) = ck_{\mu}u_s^{\mu}$. Ця формула інваріантна, її можна використовувати у будь-якій системі.

Далі усі величини розглянемо у системі спостерігача, де його 4-швидкість $u_{\text{obs}}^{\mu}=\{1,0,0,0\}$, а виміряна ним частота $\omega_{\text{obs}}=ck_{\mu}u_{\text{obs}}^{\mu}=ck_0=ck^0$. Напрям поширення хвиль, що реєструє спостерігач, є протилежним напрямку від спостерігача на джерело випромінювання $\vect{n}=-\vect{k}/|\vect{k}|=-\vect{k}/k^0$.

Після переходу в систему спостерігача маємо
\begin{equation*}
    \frac{\omega_{\text{obs}}}{\omega_0}
    = \frac{k_{\mu}u_{\text{obs}}^{\mu}}{k_{\mu}u_0^{\mu}}
    = \frac{k_0}{k_{\mu}u_0^{\mu}}
    = \frac{1}{u_0^0(1+\vect{n}\vect{v}/c)}
    = \frac{\sqrt{1-\vect{v}^2/c^2}}{1+\vect{n}\vect{v}/c},
\end{equation*}
де $\vect{v}$ --- швидкість джерела відносно спостерігача. За нерелятивістських швидкостей корінь у чисельнику можна наближено замінити на одиницю.

Нехай спостерігач та джерело знаходяться весь час на одній прямій. Коли джерело віддаляється від спостерігача, то
\begin{equation*}
    \frac{\omega_{\text{obs}}}{\omega_0}
    = \left(\frac{1-v/c}{1+v/c}\right)^{1/2} < 1, \qquad v=|\vect{v}|.
\end{equation*}
Коли, навпаки, джерело наближається до спостерігача, маємо
\begin{equation*}
    \frac{\omega_{\text{obs}}}{\omega_0}
    = \left(\frac{1+v/c}{1-v/c}\right)^{1/2} > 1.
\end{equation*}

%% --------------------------------------------------------
\section{Вектор-потенціал електромагнітного поля у випадку ізольованої системи}
%% --------------------------------------------------------

У цьому розділі отримаємо явний розв'язок рівнянь електромагнітного поля за допомогою \emph{4-вектора потенціалу\index{Чотиривектор!потенціал}}, існування якого забезпечене рівняннями~\eqref{eq:Maxwell_4D_2}: внаслідок цих рівнянь можна ввести коваріантне векторне поле $A_{\mu}$ (4-потенціал), таке, що
\begin{equation}\label{eq:4potential_def}
    F_{\mu\nu} = \partial_{\mu}A_{\nu} - \partial_{\nu}A_{\mu}.
\end{equation}
Дійсно, підстановка~\eqref{eq:4potential_def} в рівняння~\eqref{eq:Maxwell_4D_2} дає тотожний нуль. Можна показати й обернене, що з рівняння~\eqref{eq:Maxwell_4D_2} дійсно випливає існування деякого поля $A_{\mu}$, що дає~\eqref{eq:4potential_def}. Однак ми не доводимо це твердження, оскільки ми отримаємо $A_{\mu}$ в явному вигляді, з якого буде видно, що це дійсно коваріантний вектор, який реалізує розв'язок~\eqref{eq:Maxwell_4D_1}, що відповідає умовам ізольованої системи зарядів та струмів; при цьому~\eqref{eq:Maxwell_4D_2}, як було зазначено, виконуються тотожно.

Очевидно, фізичні поля, що їх описує тензор $F_{\mu\nu}$ згідно до~\eqref{eq:4potential_def}, залишаться незмінними при перетворенні виду $A_{\mu}\to A_{\mu}+\partial_{\mu}\chi$, де $\chi$ --- довільна функція. Це перетворення 4-потенціалу називають \emph{калібрувальним (або градієнтним)\index{Калібрувальне перетворення}}. Завдяки калібрувальній інваріантності можна накласти \emph{умову Лоренца\index{Умова Лоренца}}\footnote{Якщо б ця умова не виконувалася, можна підібрати функцію $\chi$ так, аби забезпечити~\eqref{eq:Lorenz_gauge} для нового 4-потенціалу.}:
\begin{equation}\label{eq:Lorenz_gauge}
    \partial_{\nu}A^{\nu} = 0.
\end{equation}
Цю умову також далі перевіримо явно, коли отримаємо потрібний розв'язок.

За умови~\eqref{eq:Lorenz_gauge} підстановка~\eqref{eq:4potential_def} у рівняння Максвелла~\eqref{eq:Maxwell_4D_1} дає
\begin{equation}\label{eq:wave_eq_A}
    \mdlgwhtsquare\ A^{\nu} = \frac{4\pi}{c}j^{\nu},
\end{equation}
де $\mdlgwhtsquare \equiv \partial_{\nu}\partial^{\nu} = \eta^{\mu\nu}\partial_{\mu}\partial_{\nu}$ --- оператор Даламбера (хвильовий оператор)\index{Оператор Даламбера}.

Розв'язок~\eqref{eq:wave_eq_A}, що задовольняє також~\eqref{eq:Lorenz_gauge}, породжує розв'язок рівнянь Максвелла для тензора електромагнітного поля за формулою~\eqref{eq:4potential_def}. Далі ми знов звернемося до постановки задачі, що відповідає \emph{ізольованій системі за відсутності зовнішнього випромінювання}. Розглянемо розв'язок~\eqref{eq:wave_eq_A} починаючи з нескінченного минулого, припускаючи, що 4-вектор $j^{\nu}(x)\equiv j^{\nu}(x^0,\vect{x})$ відмінний від нуля лише в обмеженій області просторових координат $\vect{x}$, а будь-які джерела на нескінченності відсутні. Як відомо з курсу математичної фізики, цим умовам відповідає розв'язок рівняння~\eqref{eq:wave_eq_A} у вигляді запізнюючих потенціалів\index{Запізнюючі потенціали}, що можна подати за допомогою фундаментального розв'язку $D_{\text{ret}}(x)$ оператора Даламбера:
\begin{equation*}
    D_{\text{ret}}(x) = \frac{1}{2\pi}\,\theta(x^0)\,\delta[(x^0)^2 - \vect{x}^2],
    \qquad x\equiv\{x^{\alpha}\}\equiv(x^0,\vect{x}),
\end{equation*}
де $\theta(t)\equiv\{0,\,x<0;\;1,\,x>0\}$ --- функція Хевісайда, $\delta$ --- функція Дірака; цей розв'язок, за визначенням, задовольняє рівнянню
\begin{equation*}
    \mdlgwhtsquare\ D_{\text{ret}}(x) = \delta^4(x)
    \equiv \delta(x^0)\,\delta(x^1)\,\delta(x^2)\,\delta(x^3).
\end{equation*}
Відповідний розв'язок рівняння~\eqref{eq:wave_eq_A} має вид згортки\footnote{Згортка функцій (відрізняти від згортки тензорів!) визначається за формулою
$\int f(x-x')g(x')\,d^4x'$; символ <<$*$>> у~\eqref{eq:retarded_potentials} означає саме її.}:
\begin{equation}\label{eq:retarded_potentials}
    A^{\nu} = \frac{4\pi}{c}\,D_{\text{ret}}*j^{\nu}.
\end{equation}
Формально це можна записати так:
\begin{align*}
    A^{\nu}(x) &= \frac{2}{c}\int d^4x\;
    \theta(x^0-x'^0)\,\delta[(x^0-x'^0)^2 - (\vect{x}-\vect{x}')^2]\,
    j^{\nu}(x') \\
    &= \frac{1}{c}\int d^3\vect{x}'\;
    \frac{j^{\nu}(x_{\text{ret}}^0,\vect{x}')}{|\vect{x}-\vect{x}'|},
    \qquad x_{\text{ret}}^0 \equiv x^0 - |\vect{x}-\vect{x}'|.
\end{align*}

Розв'язок описує електромагнітне поле обмеженої системи з 4-вектором густини струму $j^{\nu}$. Функція $D_{\text{ret}}(x)$ є скаляром відносно власних перетворень Лоренца, оскільки тут аргумент $\delta$-функції залежить від квадрата інтервалу, а в $\theta(x^0)$ аргумент не змінює знак при цих перетвореннях. Тому з~\eqref{eq:retarded_potentials} випливає, що $A^{\nu}$ \textbf{є чотиривектором відносно перетворень власної групи Лоренца\index{Група Лоренца}, оскільки} $j^{\nu}$ \textbf{--- чотиривектор}.

Формула~\eqref{eq:retarded_potentials}, що явно гарантує властивості $A^{\nu}$ як чотири-вектора, дозволяє прямо обчислити за допомогою~\eqref{eq:4potential_def} закон перетворення компонент $F_{\mu\nu}$ при переході в іншу систему відліку. Це дає змогу встановити, що $F_{\mu\nu}$ \textbf{--- двічі коваріантний тензор}.

Користуючись рівнянням~\eqref{eq:retarded_potentials}, перевіримо також калібрувальну умову~\eqref{eq:Lorenz_gauge}. Маємо
\begin{align*}
    \partial_{\nu}A^{\nu}
    &= \frac{4\pi}{c}\,\partial_{\nu}\{D_{\text{ret}}*j^{\nu}\}
    \equiv \frac{4\pi}{c}\int\frac{\partial}{\partial x^{\nu}}D(x-x')\,j^{\nu}(x')\,d^4x' \\
    &= \frac{4\pi}{c}\int D(x-x')\,\frac{\partial}{\partial x'^{\nu}}j^{\nu}(x')\,d^4x'
    \equiv \frac{4\pi}{c}\,D_{\text{ret}}*\partial_{\nu}j^{\nu} = 0,
\end{align*}
внаслідок закону збереження заряду~\eqref{eq:charge_conservation}. За умови Лоренца з рівнянь~\eqref{eq:Lorenz_gauge}, \eqref{eq:wave_eq_A} випливає~\eqref{eq:Maxwell_4D_1}.

Таким чином, тензор $F_{\mu\nu}$, обчислений за формулами~\eqref{eq:4potential_def}, \eqref{eq:retarded_potentials}, задовольняє рівнянням~\eqref{eq:Maxwell_4D_1}, \eqref{eq:Maxwell_4D_2}.

%=========================================================
\begin{problem}
Покладаючи $\{A^{\mu}\}=\{A^0,\vect{A}\}$ (контраваріантний вектор! Звертайте увагу на положення індексів!), де $A^0=\varphi$ --- скалярний та $\vect{A}$ --- векторний потенціали електромагнітного поля класичної електродинаміки відповідно, запишіть тривимірну форму рівнянь: (а)~$\partial_{\nu}A^{\nu}=0$; (б)~$\mdlgwhtsquare A^{\nu}=(4\pi/c)j^{\nu}$.
\end{problem}